\documentclass[aspectratio=169]{beamer}
\usetheme{Madrid}
\usecolortheme{default}
\usepackage[utf8]{inputenc}
\usepackage[spanish]{babel}
\usepackage{graphicx}
\usepackage{hyperref}

\title{Normas APA - Séptima Edición}
\subtitle{Estilo de Redacción Académica}
\author{Edison Achalma}
\institute{Universidad Nacional de San Cristóbal de Huamanga}
\date{\today}

\begin{document}

\frame{\titlepage}

\begin{frame}{Tabla de Contenido}
\tableofcontents
\end{frame}

\section{Principios de Redacción}

\begin{frame}{Características del Estilo APA}
\begin{block}{Objetivos}
\begin{itemize}
    \item Claridad en la comunicación
    \item Precisión en la expresión
    \item Objetividad científica
    \item Respeto y equidad
\end{itemize}
\end{block}

% Imagen: Diagrama con los 4 pilares del estilo APA
\begin{center}
\includegraphics[width=0.6\textwidth]{images/pilares_estilo_apa.png}
\end{center}
\end{frame}

\begin{frame}{Claridad y Concisión}
\textbf{Principio fundamental:}
\begin{quote}
Diga solo lo que necesita decir
\end{quote}

% Imagen: Ejemplo comparativo de texto verboso vs texto conciso
\begin{center}
\includegraphics[width=0.8\textwidth]{images/claridad_concision.png}
\end{center}
\end{frame}

\section{Tono y Voz}

\begin{frame}{Tono Formal}
\begin{columns}
\column{0.5\textwidth}
\textbf{Características:}
\begin{itemize}
    \item Lenguaje académico
    \item Evitar coloquialismos
    \item Sin jerga innecesaria
    \item Profesional y respetuoso
\end{itemize}

\column{0.5\textwidth}
% Imagen: Ejemplos de tono informal vs formal
\includegraphics[width=\textwidth]{images/tono_formal_informal.png}
\end{columns}
\end{frame}

\begin{frame}{Voz Activa vs Voz Pasiva}
% Imagen: Comparación de oraciones en voz activa vs pasiva con indicación de cuál usar
\begin{center}
\includegraphics[width=0.85\textwidth]{images/voz_activa_pasiva.png}
\end{center}

\begin{block}{Recomendación}
Preferir la \textbf{voz activa} para mayor claridad
\end{block}
\end{frame}

\begin{frame}{Cuándo Usar Voz Pasiva}
\textbf{La voz pasiva es apropiada cuando:}
\begin{itemize}
    \item El agente no es importante
    \item El agente es desconocido
    \item Se quiere enfatizar la acción
\end{itemize}

\vspace{0.5cm}
% Imagen: Ejemplos de uso apropiado de voz pasiva en contexto científico
\begin{center}
\includegraphics[width=0.75\textwidth]{images/ejemplos_voz_pasiva.png}
\end{center}
\end{frame}

\section{Persona Gramatical}

\begin{frame}{Primera, Segunda o Tercera Persona}
% Imagen: Tabla mostrando cuándo usar cada persona gramatical en APA
\begin{center}
\includegraphics[width=0.9\textwidth]{images/persona_gramatical.png}
\end{center}
\end{frame}

\begin{frame}{Uso de Primera Persona}
\begin{block}{APA 7 permite usar "yo" o "nosotros"}
Especialmente para:
\begin{itemize}
    \item Describir los pasos de la investigación
    \item Expresar conclusiones propias
    \item Diferenciar el trabajo propio del de otros
\end{itemize}
\end{block}

% Imagen: Ejemplos correctos de uso de primera persona
\begin{center}
\includegraphics[width=0.7\textwidth]{images/ejemplos_primera_persona.png}
\end{center}
\end{frame}

\section{Lenguaje Inclusivo}

\begin{frame}{Lenguaje Libre de Sesgos}
\textbf{APA promueve lenguaje que:}
\begin{itemize}
    \item Respete la diversidad
    \item Evite estereotipos
    \item Sea inclusivo
    \item No discrimine
\end{itemize}

\vspace{0.5cm}
% Imagen: Principios del lenguaje inclusivo según APA
\begin{center}
\includegraphics[width=0.65\textwidth]{images/lenguaje_inclusivo.png}
\end{center}
\end{frame}

\begin{frame}{Identidad de Género}
% Imagen: Tabla con ejemplos de lenguaje respetuoso para identidad de género
\begin{center}
\includegraphics[width=0.85\textwidth]{images/identidad_genero.png}
\end{center}

\begin{alertblock}{Importante}
Usar los términos que las personas utilizan para describirse a sí mismas
\end{alertblock}
\end{frame}

\begin{frame}{Evitar Sesgos por Edad}
% Imagen: Ejemplos de términos con sesgo vs sin sesgo relacionados con edad
\begin{center}
\includegraphics[width=0.8\textwidth]{images/sesgo_edad.png}
\end{center}
\end{frame}

\begin{frame}{Evitar Sesgos Raciales y Étnicos}
% Imagen: Guía de términos apropiados para grupos raciales y étnicos
\begin{center}
\includegraphics[width=0.85\textwidth]{images/sesgo_racial_etnico.png}
\end{center}
\end{frame}

\begin{frame}{Discapacidad}
\begin{block}{Principio}
Poner a la persona primero, no la condición
\end{block}

% Imagen: Comparación de lenguaje centrado en la persona vs en la discapacidad
\begin{center}
\includegraphics[width=0.8\textwidth]{images/lenguaje_discapacidad.png}
\end{center}
\end{frame}

\section{Números y Estadísticas}

\begin{frame}{Cuándo Usar Números vs Palabras}
\begin{block}{Regla General}
\begin{itemize}
    \item Números 10 y mayores: usar cifras (15, 234, 1,500)
    \item Números menores a 10: usar palabras (tres, cinco, nueve)
\end{itemize}
\end{block}

% Imagen: Tabla de reglas para números (con excepciones)
\begin{center}
\includegraphics[width=0.75\textwidth]{images/reglas_numeros.png}
\end{center}
\end{frame}

\begin{frame}{Excepciones para Números}
% Imagen: Lista de excepciones (edad, porcentajes, unidades de medida, etc.)
\begin{center}
\includegraphics[width=0.85\textwidth]{images/excepciones_numeros.png}
\end{center}
\end{frame}

\begin{frame}{Presentación de Estadísticas}
% Imagen: Ejemplos de cómo reportar estadísticas (M, SD, p, r, etc.)
\begin{center}
\includegraphics[width=0.9\textwidth]{images/presentacion_estadisticas.png}
\end{center}
\end{frame}

\section{Uso de Cursivas}

\begin{frame}{Cuándo Usar Cursivas}
\textbf{Se usa cursiva para:}
\begin{itemize}
    \item Títulos de libros, revistas, películas
    \item Introducción de términos técnicos
    \item Símbolos estadísticos (variables)
    \item Palabras en otros idiomas
    \item Énfasis (usar con moderación)
\end{itemize}
\end{frame}

\begin{frame}{Ejemplos de Uso de Cursivas}
% Imagen: Ejemplos prácticos de cada caso de uso de cursivas
\begin{center}
\includegraphics[width=0.85\textwidth]{images/ejemplos_cursivas.png}
\end{center}
\end{frame}

\section{Uso de Mayúsculas}

\begin{frame}{Reglas de Mayúsculas en APA}
% Imagen: Tabla de cuándo usar mayúsculas en títulos, nombres propios, etc.
\begin{center}
\includegraphics[width=0.9\textwidth]{images/reglas_mayusculas.png}
\end{center}
\end{frame}

\begin{frame}{Mayúsculas en Títulos}
\begin{block}{En español}
Solo la primera palabra y nombres propios
\end{block}

\begin{block}{En inglés}
Palabras principales (sustantivos, verbos, adjetivos)
\end{block}

% Imagen: Comparación de títulos en español vs inglés
\begin{center}
\includegraphics[width=0.75\textwidth]{images/mayusculas_titulos.png}
\end{center}
\end{frame}

\section{Párrafos y Transiciones}

\begin{frame}{Estructura de Párrafos}
\textbf{Un buen párrafo académico:}
\begin{itemize}
    \item Desarrolla una idea principal
    \item Tiene oración temática
    \item Incluye evidencia o ejemplos
    \item Concluye o conecta con el siguiente
\end{itemize}

% Imagen: Anatomía de un párrafo académico
\begin{center}
\includegraphics[width=0.7\textwidth]{images/estructura_parrafo.png}
\end{center}
\end{frame}

\begin{frame}{Transiciones entre Párrafos}
% Imagen: Lista de conectores y transiciones útiles organizados por función
\begin{center}
\includegraphics[width=0.85\textwidth]{images/conectores_transiciones.png}
\end{center}
\end{frame}

\section{Errores Comunes}

\begin{frame}{Errores Frecuentes de Estilo}
% Imagen: Lista de los 10 errores más comunes con ejemplos
\begin{center}
\includegraphics[width=0.9\textwidth]{images/errores_comunes_estilo.png}
\end{center}
\end{frame}

\begin{frame}{Redundancias a Evitar}
% Imagen: Ejemplos de redundancias comunes y cómo corregirlas
\begin{center}
\includegraphics[width=0.85\textwidth]{images/redundancias_evitar.png}
\end{center}
\end{frame}

\begin{frame}{Lista de Verificación de Estilo}
% Imagen: Checklist para revisar el estilo de redacción
\begin{center}
\includegraphics[width=0.8\textwidth]{images/checklist_estilo.png}
\end{center}
\end{frame}

\begin{frame}[standout]
\Huge ¿Preguntas?
\end{frame}

\begin{frame}{Próxima Clase}
\begin{center}
\Large \textbf{Citas en el Texto}
\vspace{1cm}

Veremos:
\begin{itemize}
    \item Citas textuales cortas y largas
    \item Parafraseo
    \item Citas narrativas vs parentéticas
    \item Reglas según número de autores
\end{itemize}
\end{center}
\end{frame}

\end{document}
